\section{Revisão da Literatura}

% ---------------------------------------------------------------- 6
\begin{frame}{PINNs para EDPs da hidrodinâmica}
  Uma Physics-Informed Neural Network (PINN) \parencite{raissi2019} adiciona o resíduo
  da EDP na função objetivo, construindo uma rede que aproxima a solução da equação/respeita a lei física.
  \vspace{4pt}
  \begin{center}
    \scriptsize
    \renewcommand{\arraystretch}{1.3}
    \begin{tabular}{@{}l l p{7.0cm}@{}}
      \toprule
      \textbf{Trabalho} & \textbf{Equação} & \textbf{Forma resolvida} \\
      \midrule
      \textcite{bihlo2022}  & \parbox[t]{2.7cm}{\raggedright Águas rasas na esfera} &
        $(\eta - h)_t + \nabla\!\cdot\bigl((\eta - h)\,\mathbf{u}\bigr) = 0$,\newline
        $\mathbf{u}_t + (\mathbf{u}\cdot\nabla)\mathbf{u} + f\,\mathbf{k}\times\mathbf{u} + g\,\nabla \eta = 0$ \\
      \textcite{jagtap2022} & \parbox[t]{2.7cm}{\raggedright Serre--Green--Naghdi (1D)} &
        $(h + \eta)_t + \bigl((h + \eta)\,u\bigr)_x = 0$,\newline
        $u_t + g\,\eta_x + u\,u_x
          - \tfrac{1}{h + \eta}\bigl[(h + \eta)^2\bigl(\tfrac{1}{3}P + \tfrac{1}{2}Q\bigr)\bigr]_x$\newline
        \hspace*{1.2em}${} + \bigl(\tfrac{1}{2}P + Q\bigr) h_x = 0$ \\
      \textcite{zheng2023}  & \parbox[t]{2.7cm}{\raggedright Boussinesq $(2+1)$-dimensional} & --- \\
      \textcite{wang2024}   & \parbox[t]{2.7cm}{\raggedright Boussinesq e Camassa--Holm} &
        $\eta_{tt} - \eta_{xx} - 3\,(\eta^2)_{xx} - \eta_{xxxx} = 0$,\newline
        $u_t + 2\kappa\,u_x - u_{xxt} + 3\,u\,u_x - 2\,u_x u_{xx} - u\,u_{xxx} = 0$ \\
      \bottomrule
    \end{tabular}
  \end{center}
  \vspace{2pt}
  \textcolor{AmurmapleRed}{Note que não há trabalhos de SciML para o sistema de Boussinesq na forma paramétrica de \textcite{martins2023}, que é o foco deste trabalho.}
\end{frame}

% ---------------------------------------------------------------- 7
\begin{frame}{Aprendizado de operadores em fluidos}
  Outro braço de métodos dentro de SciML, buscando aproximar o operador que leva o
  parâmetro da EDP à solução (operador entre espaços de função).
  \vspace{4pt}
  \begin{center}
    \small
    \begin{tabular}{@{}ll@{}}
      \toprule
      \textbf{Trabalho} & \textbf{Contribuição} \\
      \midrule
      \textcite{lu2021deeponet} & DeepONet: primeira arquitetura profunda para operadores \\
      \textcite{li2021}         & FNO: kernel em Fourier, independente da discretização \\
      \textcite{gupta2021}      & bases alternativas, como wavelets \\
      \textcite{li2023pino}     & PINO: resíduo da EDP no operador \\
      \bottomrule
    \end{tabular}
  \end{center}
  \vspace{2pt}
  \begin{itemize}[<+->]
    \item Escala geofísica: oceano regional e previsão atmosférica global
          \parencite{bonev2023, uchida2025, yu2025};
    \item Independência da discretização (zero-shot superresolution), possibilidade de avaliar em uma resolução maior do que a treinada \parencite{li2021, li2023pino}.
    \item \textcolor{AmurmapleRed}{Também não há trabalhos aplicando para o sistema de Boussinesq}
  \end{itemize}
\end{frame}

% ---------------------------------------------------------------- 8
\begin{frame}{Viés espectral e o Neural Tangent Kernel}
  Uma característica presente ao utilizar métodos de SciML para a solução de EDPs de evolução, é o \textbf{viés espectral}, que consiste em no fato do método priorizar aproximar as baixas frequências da solução e necessitar de muito mais iterações para as altas.

  \begin{itemize}[<+->]
    \item O viés espectral foi inicialmente relatado por \textcite{rahaman2019, xu2019frequency}, apresentado para PINNs por \textcite{wang2022ntk} e para operadores neurais \textcite{li2023pino, khodakarami2026};
    \item Explicações teóricas que encontramos são restritas a MLPs específicas, como uma rede de duas camadas com ReLU \parencite{arora2019, cao2019, basri2019}.
  \end{itemize}
  \begin{information}[O que fizemos]
    Usamos o framework de Neural Tangent Kernel (NTK) \parencite{jacot2018} para isolar o viés espectral para uma MLP simples, encontrando estimativas do treinamento (válidas sob certas hipóteses)
    \parencite{chizat2019}.
  \end{information}
\end{frame}
